\documentclass[10pt]{article}
\usepackage[T1]{fontenc}
\usepackage[utf8]{inputenc}
\usepackage[french]{babel}
\usepackage{amsmath,amssymb,mathrsfs,mathpazo}
\usepackage{geometry}
\geometry{a4paper,margin=2cm}
\usepackage{graphicx,array,tabularx,tabularray,colortbl,multirow,multicol,booktabs,enumitem,pifont,ulem,fancybox,diagbox,cancel}
\usepackage{pstricks,pstricks-add,pst-plot,pst-tree,pst-node,pst-text,pst-func,pst-math,pst-bspline,pst-3dplot,pst-eucl,pst-solides3d}
\let\pscalculate\undefined
\usepackage{pst-bezier}
\usepackage{tikz,pgfplots,tkz-tab}
\usetikzlibrary{arrows,arrows.meta,patterns,calc,positioning,trees,decorations.markings,shapes.geometric,shapes.misc,shapes.arrows,intersections,decorations.pathmorphing,decorations.pathreplacing,backgrounds,shadows,babel}
\pgfplotsset{compat=1.18}
\usepgfplotslibrary{fillbetween}
\usepackage[np]{numprint}
\usepackage{listings,makecell,pas-tableur,twoopt}
\let\barre\undefined
\usepackage[pstricks]{tablvar}
\definecolor{darkgreen}{RGB}{0,100,0}
\providecommand{\red}{\color{red}}
\providecommand{\blue}{\color{blue}}
\providecommand{\green}{\color{green}}
\providecommand{\black}{\color{black}}
\providecommand{\white}{\color{white}}
\providecommand{\gray}{\color{gray}}
\providecommand{\cyan}{\color{cyan}}
\providecommand{\magenta}{\color{magenta}}
\providecommand{\yellow}{\color{yellow}}
\providecommand{\violet}{\color{violet}}
\providecommand{\orange}{\color{orange}}
\providecommand{\euro}{\texteuro}
\providecommand{\EUR}{\texteuro}
\providecommand{\ieme}{\textsuperscript{e}}
\providecommand{\degres}{\ensuremath{^\circ}}
\providecommand{\degree}{\ensuremath{^\circ}}
\providecommand{\vv}[1]{\overrightarrow{#1}}
\providecommand{\Coord}[1]{#1}
\providecommand{\decosix}{}
\providecommand{\starredbullet}{\ensuremath{\bullet}}
\pagestyle{empty}
\setlength{\parindent}{0pt}
\usepackage{adjustbox}\DeclareUnicodeCharacter{2236}{:}\DeclareUnicodeCharacter{2190}{\ensuremath{\leftarrow}}
\providecommand{\R}{\mathbb{R}}
\providecommand{\N}{\mathbb{N}}
\providecommand{\Z}{\mathbb{Z}}
\providecommand{\Q}{\mathbb{Q}}
\providecommand{\D}{\mathbb{D}}
\providecommand{\C}{\mathbb{C}}
\providecommand{\e}{{\mathrm{e}}}
\providecommand{\vect}[1]{\overrightarrow{#1}}
\providecommand{\vectt}[1]{\overrightarrow{\mathrm{#1}}}
\providecommand{\vvt}[1]{\overrightarrow{\mathrm{#1}}}
\providecommand{\np}[2][]{{#2~#1}}
\providecommand{\vv}[1]{{\overrightarrow{#1}}}
\renewcommand{\d}{{\mathrm{d}}}
\providecommand{\degres}{{^\circ}}
\providecommand{\degre}{{^\circ}}
\providecommand{\blue}{\color{blue}}
\providecommand{\red}{\color{red}}
\providecommand{\enskip}{\quad}
\providecommand{\llbracket}{\lbrack\!\lbrack}
\providecommand{\rrbracket}{\rbrack\!\rbrack}
\providecommand{\euro}{{\text{€}}}
\providecommand{\no}{{\text{n°}}}
\providecommand{\interval}[2]{{[#1;#2]}}
\providecommand{\intervFO}[2]{{[#1;#2[}}
\providecommand{\dotfill}{\ldots}

\providecommand{\ds}{\displaystyle}
\providecommand{\textding}[1]{#1}
\providecommand{\Oij}{\(\mathrm{O};\vec{\imath},\vec{\jmath}\)}
\providecommand{\Oijk}{\(\mathrm{O};\vec{\imath},\vec{\jmath},\vec{k}\)}
\providecommand{\Ouv}{\(\mathrm{O};\vec{u},\vec{v}\)}
\providecommand{\intervFF}[2]{\ensuremath{[#1;#2]}}
\renewcommand{\euro}{\ifmmode\text{\texteuro}\else\texteuro\fi}
\DeclareUnicodeCharacter{00B0}{\ifmmode^{\circ}\else\textdegree\fi}
\DeclareUnicodeCharacter{20AC}{\ifmmode\text{\texteuro}\else\texteuro\fi}
% Énoncés : Éducation nationale. Transcription et corrigés : APMEP, Denis Vergès et contributeurs. Découpage et classement : CoopMaths, Sébastien Lozano, Éric Elter et contributeurs. Adaptation de la présentation : Maths Explorer.
% CC BY-SA : https://coopmaths.fr/www/about/
% https://coopmaths.fr/alea/static/bac/2023/tex/bac_2023_03_sujet2_asie_2.tex
\begin{document}
\noindent\textit{APMEP / CoopMaths. CC BY-SA. Présentation adaptée par Maths Explorer.}\par\medskip
\medskip

On considère la fonction $f$ définie sur $\R$ par 

\[f(x) = \ln \left(\e^{2x}  - \e^x  + 1\right).\]

On note $\mathcal{C}_f$ sa courbe représentative représentée ci-dessous.

\begin{center}
\psset{unit=1.35cm,arrowsize=2pt 3,comma=true}
\begin{pspicture*}(-3.6,-2)(3,3)
\psaxes[linewidth=1.25pt,Dx=0.5,Dy=0.5,labelFontSize=\scriptstyle]{->}(0,0)(-3.5,-1.95)(3,3)
\psgrid[gridlabels=0pt,subgriddiv=2]
\psplot[plotpoints=2000,linewidth=1.25pt,linecolor=blue]{-3.5}{2}{2.71828 x 2 mul exp 2.71828 x exp sub 1 add ln} \uput[dr](1.4,2.5){\blue $\mathcal{C}_f$}
\end{pspicture*}
\end{center}

Un élève formule les conjectures suivantes à partir de cette représentation graphique :

\smallskip

\begin{tabularx}{\linewidth}{|X|}\hline
\textbf{1.~~}L'équation $f(x) = 2$ semble admettre au moins une solution.\\
\textbf{2.~~}Le plus grand intervalle sur lequel la fonction $f$ semble être croissante est $[- 0,5~;~+\infty[$.\\
\textbf{3.~~}L'équation de la tangente au point d'abscisse $x = 0$ semble être: $y = 1,5x$.\\ \hline
\end{tabularx}

\medskip

Le but de cet exercice est de valider ou rejeter les conjectures concernant la fonction $f$.

\medskip

\textbf{Partie A : Étude d'une fonction auxiliaire}

\medskip

On définit sur $\R$ la fonction $g$ définie par 
\[g(x) = \e^{2x} - \e^x + 1.\]

\begin{enumerate}
\item Déterminer $\displaystyle\lim_{x \to - \infty} g(x)$.
\item Montrer que $\displaystyle\lim_{x \to + \infty} g(x) = + \infty$.
\item Montrer que $g'(x) = \e^x\left(2\e^x - 1\right)$ pour tout $x \in  \R$.
\item Étudier le sens de variation de la fonction $g$ sur $\R$.

Dresser le tableau des variations de la fonction $g$ en y faisant figurer la valeur exacte des extremums s'il y en a, ainsi que les limites de $g$ en $-\infty$ et $+ \infty$.
\item En déduire le signe de $g$ sur $\R$.
\item Sans en mener nécessairement les calculs, expliquer comment on pourrait établir le résultat
de la question 5 en posant $X = \e^x$.
\end{enumerate}

\medskip

\textbf{Partie B}

\medskip

\begin{enumerate}
\item Justifier que la fonction $f$ est bien définie sur $\R$.
\item La fonction dérivée de la fonction $f$ est notée $f'$.

Justifier que $f'(x) = \dfrac{g'(x)}{g(x)}$ pour tout $x \in \R$.
\item Déterminer une équation de la tangente à la courbe au point d'abscisse $0$.
\item Montrer que la fonction $f$ est strictement croissante sur $[-\ln (2)~;~ +\infty[$.
\item Montrer que l'équation $f(x) = 2$ admet une unique solution $\alpha$ sur $[-\ln (2)~;~ +\infty[$ et déterminer une valeur approchée de $\alpha$ à $10^{-2}$ près.
\end{enumerate}

\medskip

\textbf{Partie C}

\medskip

À l'aide des résultats de la partie B, indiquer, pour chaque conjecture de l'élève, si elle est vraie ou fausse. Justifier.

\medskip

\end{document}
